(a) Fix $x$ with $\varphi(x)=r$. Lift a basis of the fiber to local sections of $\mathcal F$. By Nakayama's lemma, their cokernel has zero stalk at $x$. Its support is closed, so these $r$ sections generate on a neighborhood of $x$. Thus $\varphi\le r$ there. Consequently $\{x:\varphi(x)<n\}$ is open for every $n$.

(b) On an open set where $\mathcal F\cong\mathcal O_X^r$, the function $\varphi$ equals $r$. Hence it is locally constant, and connectedness makes it constant.

(c) Suppose $\varphi=r$. As in (a), near any point there is a surjection $\mathcal O_U^r\to\mathcal F|_U$. Take $U=\operatorname{Spec}A$. For any relation $(a_1,\ldots,a_r)$ among these generators and any prime $\mathfrak p$, the induced map $k(\mathfrak p)^r\to\mathcal F\otimes k(\mathfrak p)$ is an isomorphism. Hence every $a_i$ lies in every prime of $A$. Since $A$ is reduced, all $a_i=0$. The surjection is therefore an isomorphism.
